\documentclass[11pt,a4paper]{article}

\usepackage[utf8]{inputenc}
\usepackage[T1]{fontenc}
\usepackage{amsmath,amssymb,amsthm}
\usepackage{hyperref}
\usepackage{graphicx}
\usepackage{booktabs}
\usepackage[margin=2.5cm]{geometry}
\usepackage{natbib}
\usepackage{algorithm}
\usepackage{algpseudocode}

\newtheorem{definition}{Definition}
\newtheorem{proposition}{Proposition}

\title{LLM-Based Novelty Detection in Academic Text:\\An Information-Theoretic Approach}
\author{Simon McCallum}
\date{\today}

\begin{document}
\maketitle

\begin{abstract}
We present a method for detecting novel content within academic documents using
Large Language Models (LLMs) as semantic compressors combined with FAISS-based
similarity search. Each text chunk is processed by an LLM to generate a concise
semantic prompt, which is then embedded in a vector space. Novelty is quantified
as the inverse of average cosine similarity to nearest neighbors, capturing
semantic uniqueness rather than lexical distinctiveness. We provide an
information-theoretic justification grounded in Kolmogorov complexity and discuss
the method's advantages over traditional approaches.
\end{abstract}

\section{Introduction}

Detecting which parts of a document contain genuinely novel contributions versus
standard or repetitive content is a fundamental challenge in academic writing,
peer review, and content analysis. Traditional approaches based on $n$-gram
overlap, TF-IDF, or lexical matching fail to capture semantic novelty---they
cannot distinguish between a paraphrased common idea and a truly original
contribution expressed in familiar terms.

We propose a system that leverages the semantic understanding of LLMs to detect
intra-document novelty. The key insight is that an LLM, when asked to summarize
or ``compress'' a text passage, produces output whose information content
reflects the novelty of the input relative to the model's training
distribution.

\section{Information-Theoretic Foundation}

\subsection{Kolmogorov Complexity}

The Kolmogorov complexity $K(x)$ of a string $x$ is the length of the shortest
program that produces $x$ on a universal Turing machine
\citep{li2008introduction}:

\begin{definition}
$K(x) = \min\{|p| : U(p) = x\}$
\end{definition}

While $K(x)$ is uncomputable in general, practical compression algorithms
provide upper bounds. We extend this notion to semantic content, using LLMs as
approximate semantic compressors.

\subsection{Conditional Complexity and Document Context}

For a chunk $c_i$ within document $D$, the conditional Kolmogorov complexity
$K(c_i | D \setminus c_i)$ captures the information content of $c_i$ given the
rest of the document. High conditional complexity implies the chunk contains
information not derivable from the remaining text---i.e., novel content.

\begin{proposition}
If $c_i$ is semantically redundant with other parts of $D$, then
$K(c_i | D \setminus c_i) \ll K(c_i)$.
\end{proposition}

\section{LLM as Semantic Compressor}

\subsection{Prompt Generation as Lossy Compression}

Given a text chunk $c_i$ with context $(c_{i-1}, c_{i+1})$, we query an LLM
$\mathcal{L}$ to produce a semantic prompt $p_i$:

\begin{equation}
p_i = \mathcal{L}(c_i, c_{i-1}, c_{i+1})
\end{equation}

This prompt is a lossy compression of $c_i$ that preserves its essential
semantic content while discarding surface-level details. The LLM acts as an
approximate semantic compressor, where:

\begin{itemize}
\item \textbf{Common content} (literature review boilerplate, standard
methodology descriptions) produces generic prompts that are similar across
chunks, reflecting low conditional complexity.
\item \textbf{Novel content} (original analysis, unique findings) produces
specific, distinctive prompts, reflecting high conditional complexity.
\end{itemize}

\subsection{Information Preservation}

The prompt $p_i$ preserves the semantic information of $c_i$ in the sense that:

\begin{equation}
I(c_i; p_i) \approx H_{sem}(c_i)
\end{equation}

where $I(c_i; p_i)$ is the mutual information between chunk and prompt, and
$H_{sem}(c_i)$ is the semantic entropy of the chunk. The LLM's training on
large corpora provides an implicit model of ``common'' vs.\ ``uncommon''
content, making its compression behaviour sensitive to novelty.

\section{Embedding Space Semantics}

\subsection{Sentence-BERT Embeddings}

We encode each prompt $p_i$ into a dense vector $\mathbf{e}_i \in
\mathbb{R}^d$ using a pre-trained Sentence-BERT model
\citep{reimers2019sentence}:

\begin{equation}
\mathbf{e}_i = \text{SBERT}(p_i) \in \mathbb{R}^d
\end{equation}

The embedding space is structured such that semantically similar texts have
high cosine similarity:

\begin{equation}
\text{sim}(p_i, p_j) = \cos(\mathbf{e}_i, \mathbf{e}_j) = \frac{\mathbf{e}_i \cdot \mathbf{e}_j}{\|\mathbf{e}_i\| \|\mathbf{e}_j\|}
\end{equation}

After L2 normalization ($\hat{\mathbf{e}}_i = \mathbf{e}_i / \|\mathbf{e}_i\|$),
cosine similarity reduces to inner product:

\begin{equation}
\cos(\hat{\mathbf{e}}_i, \hat{\mathbf{e}}_j) = \hat{\mathbf{e}}_i \cdot \hat{\mathbf{e}}_j
\end{equation}

This property enables efficient similarity search using inner product indices.

\section{FAISS Similarity Search}

We use Facebook AI Similarity Search (FAISS) \citep{johnson2019billion} with an
\texttt{IndexFlatIP} (flat inner product) index for exact nearest neighbor
search on L2-normalized vectors. For each embedding $\hat{\mathbf{e}}_i$, we
retrieve the $k$ nearest neighbors:

\begin{equation}
\mathcal{N}_k(i) = \underset{j \neq i, |\mathcal{S}|=k}{\arg\max} \sum_{j \in \mathcal{S}} \hat{\mathbf{e}}_i \cdot \hat{\mathbf{e}}_j
\end{equation}

The computational complexity is $O(nd)$ per query for exact search with $n$
vectors of dimension $d$, which is practical for document-length texts (typically
$n < 500$ chunks).

\section{Novelty Score Formulation}

\subsection{Definition}

The novelty score for chunk $c_i$ is defined as:

\begin{equation}
\text{novelty}(c_i) = 1 - \frac{1}{k} \sum_{j \in \mathcal{N}_k(i)} \cos(\hat{\mathbf{e}}_i, \hat{\mathbf{e}}_j)
\end{equation}

\subsection{Properties}

\begin{itemize}
\item $\text{novelty}(c_i) \in [0, 1]$ (since cosine similarity of normalized
vectors is in $[-1, 1]$, though in practice prompts from the same document have
positive similarity)
\item $\text{novelty}(c_i) \to 0$ when the chunk's prompt is nearly identical to
its neighbors (highly repetitive)
\item $\text{novelty}(c_i) \to 1$ when the chunk's prompt is orthogonal to all
neighbors (completely unique)
\end{itemize}

\subsection{Score Categories}

We define four novelty categories with empirically determined thresholds:

\begin{table}[h]
\centering
\begin{tabular}{lll}
\toprule
Category & Score Range & Interpretation \\
\midrule
High Novelty & $> 0.7$ & Unique content, original analysis \\
Medium Novelty & $0.4$--$0.7$ & Somewhat unique, builds on themes \\
Low Novelty & $0.2$--$0.4$ & Common content, standard methodology \\
Very Low Novelty & $< 0.2$ & Repetitive, boilerplate \\
\bottomrule
\end{tabular}
\caption{Novelty score categories and interpretations.}
\label{tab:categories}
\end{table}

\section{Algorithm}

\begin{algorithm}
\caption{LLM-Based Novelty Detection}
\begin{algorithmic}[1]
\Require PDF document $D$, chunk size $w$, overlap $o$, neighbors $k$
\Ensure Novelty scores for each chunk
\State $T \gets \text{ExtractText}(D)$
\State $C \gets \text{Chunk}(T, w, o)$ \Comment{List of text chunks}
\For{each chunk $c_i \in C$}
    \State $p_i \gets \mathcal{L}(c_i, c_{i-1}, c_{i+1})$ \Comment{LLM semantic compression}
\EndFor
\State $E \gets \text{SBERT}(p_1, \ldots, p_n)$ \Comment{Embed all prompts}
\State $\hat{E} \gets \text{L2Normalize}(E)$
\State $\text{Index} \gets \text{FAISS.IndexFlatIP}(\hat{E})$
\For{each embedding $\hat{\mathbf{e}}_i$}
    \State $\mathcal{N}_k(i) \gets \text{Index.search}(\hat{\mathbf{e}}_i, k+1) \setminus \{i\}$
    \State $s_i \gets 1 - \frac{1}{k}\sum_{j \in \mathcal{N}_k(i)} \hat{\mathbf{e}}_i \cdot \hat{\mathbf{e}}_j$
\EndFor
\State \Return $\{(c_i, s_i)\}_{i=1}^n$
\end{algorithmic}
\end{algorithm}

\section{Advantages}

\begin{enumerate}
\item \textbf{Semantic over lexical}: Captures paraphrased redundancy that
lexical methods miss. Two chunks expressing the same idea in different words
will produce similar prompts and thus low novelty scores.

\item \textbf{Context-aware}: The LLM receives surrounding context, enabling it
to distinguish between a concept's first introduction and its later repetition.

\item \textbf{Multi-provider flexibility}: The architecture supports multiple
LLM providers (local Ollama, cloud APIs) with different cost-quality
trade-offs, making it accessible for various budgets.

\item \textbf{Efficient search}: FAISS provides sub-linear search time for
larger documents, though exact search is sufficient for typical academic papers
($< 500$ chunks).
\end{enumerate}

\section{Limitations and Future Work}

\begin{itemize}
\item \textbf{Intra-document only}: The current system measures novelty relative
to the document itself, not against external corpora or the broader literature.

\item \textbf{Threshold sensitivity}: The novelty category boundaries ($0.2$,
$0.4$, $0.7$) are empirically set and may need calibration for different
document types.

\item \textbf{LLM quality dependence}: The quality of novelty detection depends
on the LLM's ability to generate semantically meaningful prompts. Smaller models
(e.g., local Ollama) may produce less discriminative prompts than larger cloud
models.

\item \textbf{Computational cost}: Each chunk requires an LLM API call,
making analysis of large documents potentially expensive with cloud providers.
Local Ollama mitigates cost but may sacrifice quality.

\item \textbf{Future directions}: Integration with external knowledge bases for
cross-document novelty, fine-tuning embedding models for academic text,
and development of domain-specific novelty thresholds.
\end{itemize}

\bibliographystyle{plainnat}
\begin{thebibliography}{9}

\bibitem[Johnson et~al., 2019]{johnson2019billion}
Johnson, J., Douze, M., and J{\'e}gou, H. (2019).
\newblock Billion-scale similarity search with GPUs.
\newblock {\em IEEE Transactions on Big Data}, 7(3):535--547.

\bibitem[Li and Vit{\'a}nyi, 2008]{li2008introduction}
Li, M. and Vit{\'a}nyi, P. (2008).
\newblock {\em An Introduction to Kolmogorov Complexity and Its Applications}.
\newblock Springer, 3rd edition.

\bibitem[Reimers and Gurevych, 2019]{reimers2019sentence}
Reimers, N. and Gurevych, I. (2019).
\newblock Sentence-BERT: Sentence embeddings using Siamese BERT-Networks.
\newblock In {\em Proceedings of EMNLP-IJCNLP 2019}, pages 3982--3992.

\end{thebibliography}

\end{document}
